% !TEX root = ../main.tex
\chapter{培训与支持计划}

\section{用户培训方案}

\subsection{培训目标}
\begin{enumerate}
    \item 掌握系统登录和基本操作流程
    \item 理解负载预测和优化结果的含义
    \item 能够根据实际场景调整优化参数
    \item 具备初步的故障排查能力
\end{enumerate}

\subsection{培训内容大纲}

表~\ref{tab:training_outline}展示了培训课程的详细大纲。

\begin{table}[H]
    \centering
    \caption{用户培训课程大纲}
    \label{tab:training_outline}
    \begin{tabular}{clcc}
        \toprule
        \textbf{模块}                     & \textbf{内容}  & \textbf{时长} & \textbf{方式} \\
        \midrule
        \multirow{3}{*}{模块一：系统概述}
                                        & 系统功能介绍       & 15分钟        & 讲授          \\
                                        & 技术架构简介       & 10分钟        & 讲授          \\
                                        & 典型应用场景       & 10分钟        & 案例分析        \\
        \midrule
        \multirow{4}{*}{模块二：基本操作}
                                        & 账户注册与登录      & 10分钟        & 演示+实操       \\
                                        & 界面导航与功能模块    & 15分钟        & 演示          \\
                                        & 参数配置与设置      & 15分钟        & 实操          \\
                                        & 数据文件上传       & 10分钟        & 实操          \\
        \midrule
        \multirow{3}{*}{模块三：核心功能}
                                        & 运行能源优化       & 20分钟        & 演示+实操       \\
                                        & 解读优化结果       & 15分钟        & 案例分析        \\
                                        & 场景对比分析       & 15分钟        & 实操          \\
        \midrule
        \multirow{2}{*}{模块四：进阶使用}
                                        & 特征重要性分析      & 15分钟        & 讲授          \\
                                        & 模型健康度监控      & 10分钟        & 演示          \\
        \midrule
        \multirow{2}{*}{模块五：问题处理}
                                        & 常见问题与解决方案    & 15分钟        & FAQ         \\
                                        & 技术支持渠道       & 5分钟         & 讲授          \\
        \midrule
        \multicolumn{2}{r}{\textbf{总计}} & \textbf{3小时} &                           \\
        \bottomrule
    \end{tabular}
\end{table}

\subsection{培训材料}

培训将提供以下配套材料：

\begin{enumerate}
    \item \textbf{用户操作手册}：详细的图文操作指南（PDF格式）
    \item \textbf{视频教程}：每个功能模块的操作演示视频（5-10分钟/集）
    \item \textbf{快速参考卡}：常用操作的单页速查表
    \item \textbf{FAQ文档}：常见问题与解答汇编
\end{enumerate}

\section{技术支持计划}

\subsection{支持层级}

技术支持采用三级支持体系，如图~\ref{fig:support_levels}所示。

\begin{figure}[htbp]
    \centering
    \begin{tikzpicture}[
            node distance=0.8cm,
            level/.style={rectangle, draw, rounded corners, minimum width=4cm, minimum height=1cm, align=center, font=\small}
        ]
        \node[level, fill=green!30] (l1) {一级支持\\用户自助 + FAQ};
        \node[level, fill=yellow!30, below=of l1] (l2) {二级支持\\开发团队远程协助};
        \node[level, fill=red!30, below=of l2] (l3) {三级支持\\紧急问题专项处理};

        \node[right=2cm of l1, font=\footnotesize, align=left] {
            • 查阅用户手册\\
            • 检索FAQ\\
            • 社区问答
        };

        \node[right=2cm of l2, font=\footnotesize, align=left] {
            • 邮件工单\\
            • 远程诊断\\
            • 响应时间: 24h
        };

        \node[right=2cm of l3, font=\footnotesize, align=left] {
            • 系统故障\\
            • 数据丢失\\
            • 响应时间: 4h
        };

        \draw[->, thick] (l1) -- (l2);
        \draw[->, thick] (l2) -- (l3);
    \end{tikzpicture}
    \caption{三级技术支持体系}
    \label{fig:support_levels}
\end{figure}

\subsection{支持渠道}

技术支持提供多种联系渠道：在线文档(7×24即时)、电子邮件(工作日24h响应)、GitHub Issues(48h响应)和紧急热线(4h响应)。

\subsection{问题升级流程}

问题升级遵循以下流程：

\begin{enumerate}
    \item 用户提交问题描述和截图
    \item 支持人员进行问题分类和优先级评估
    \item 低优先级问题由一级支持处理
    \item 技术性问题升级至开发团队
    \item 紧急问题启动专项处理流程
    \item 问题解决后进行根因分析和知识库更新
\end{enumerate}

\section{系统维护手册}

\subsection{日常维护任务}

表~\ref{tab:maintenance_tasks}定义了系统日常维护任务。

\begin{table}[htbp]
    \centering
    \caption{日常维护任务清单}
    \label{tab:maintenance_tasks}
    \begin{tabular}{llll}
        \toprule
        \textbf{任务} & \textbf{频率} & \textbf{责任人} & \textbf{检查项}      \\
        \midrule
        服务健康检查      & 每日          & 运维           & API响应状态、错误日志      \\
        数据采集验证      & 每日          & 运维           & CAISO/Weather数据更新 \\
        存储容量监控      & 每周          & 运维           & Cloud Storage使用率  \\
        模型性能评估      & 每周          & 开发           & MAPE、R²指标         \\
        安全审计        & 每月          & 安全           & 访问日志、异常登录         \\
        依赖更新        & 每季度         & 开发           & 安全补丁、版本升级         \\
        \bottomrule
    \end{tabular}
\end{table}

\subsection{故障排查指南}

\subsubsection{常见问题诊断}

\paragraph{问题1: API响应超时}
\begin{enumerate}
    \item 检查GAE实例状态（Cloud Console）
    \item 查看请求日志中的错误信息
    \item 确认Gurobi许可证是否有效
    \item 检查外部API（CAISO/Weather）连通性
\end{enumerate}

\paragraph{问题2: 优化结果异常}
\begin{enumerate}
    \item 验证输入参数是否在合理范围
    \item 检查预测负载数据是否完整
    \item 确认电价配置是否正确
    \item 查看Gurobi求解日志
\end{enumerate}

\paragraph{问题3: 前端页面加载失败}
\begin{enumerate}
    \item 清除浏览器缓存并刷新
    \item 检查网络连接状态
    \item 验证Firebase Hosting部署状态
    \item 查看浏览器控制台错误信息
\end{enumerate}

\subsection{应急响应流程}

对于系统故障，启动以下应急响应流程：

\begin{figure}[htbp]
    \centering
    \begin{tikzpicture}[
            node distance=0.6cm,
            stepbox/.style={rectangle, draw, rounded corners, minimum width=1.8cm, minimum height=0.7cm, align=center, font=\tiny},
            arrow/.style={->, >=stealth}
        ]
        \node[stepbox, fill=red!30] (detect) {故障检测};
        \node[stepbox, fill=orange!30, right=of detect] (assess) {评估};
        \node[stepbox, fill=yellow!30, right=of assess] (notify) {通知};
        \node[stepbox, fill=blue!20, right=of notify] (fix) {修复};
        \node[stepbox, fill=green!30, right=of fix] (verify) {验证};
        \node[stepbox, fill=gray!20, right=of verify] (review) {复盘};

        \draw[arrow] (detect) -- (assess);
        \draw[arrow] (assess) -- (notify);
        \draw[arrow] (notify) -- (fix);
        \draw[arrow] (fix) -- (verify);
        \draw[arrow] (verify) -- (review);
    \end{tikzpicture}
    \caption{应急响应流程}
    \label{fig:incident_response}
\end{figure}

\subsection{运维监控仪表盘}

系统运维通过Cloud Monitoring监控以下关键指标：请求量(QPS)、延迟(P50/P95/P99)、错误率(4xx/5xx)、资源使用(CPU/内存)、定时任务执行状态。

